\section{Results}
\label{sec:results}

\subsection{Evaluation populations and computational budgets}
One shared model was trained across 22 families with 402,336 presentations per family
(8,851,392 total), using 2,794 optimizer updates and deterministic FP32 execution on eight H100s.
Each family participated in 381 updates. The complete execution took 1,377.7 seconds, including
preparation and checkpoints (Table~\ref{tab:22-training}). The principal development population
contains 1,408 TRAIN-derived requests, 64 per family, evaluated using the original checkpoint and
the developed assembly, proposal and selection procedure. All requests, including failures,
remain in the denominator. This population was used during method development and is not an
independent held-out evaluation. Independent training-seed estimates and final held-out results
remain \PendingValue{independent-seed held-out evaluation}.

Separate paired diagnostics use 352 requests (16 per family), one draw per request and 64
sampling steps. They compare raw flow endpoints, terminal argmax and the bounded D1 completion
procedure, without the later full proposal and selection procedure used for the 1,408-request
population. These populations and inference budgets are not interchangeable. Means and sample
standard deviations across independent fits are not estimable from one original fit and its
continuations. Post-hoc rejection cannot increase a proposal's per-attempt acceptance probability
(Proposition~\ref{prop:posthoc-budget}); larger candidate pools and changed selection are therefore
reported separately from learned-model comparisons.

\subsection{Shared generation across 22 assembly families}
The selected development outputs achieve exact L1 in 1,387/1,408 requests (98.5\%), with all 22
families above 90\% (Table~\ref{tab:22-l1}). The lowest family yield is 58/64 (90.6\%) for aryl
reductive amination. Exact L1 requires molecular validation, an admitted decomposition and exact
forward replay. The equal-family and pooled yields coincide because every family has 64 requests.
These counts describe the combined model and inference procedure. Reconstructing terminal argmax
from the saved logits gives 379/1,408 exact outputs (26.9\%) on the first draw and 1,910/7,040
(27.1\%) across all five draws, retaining every trajectory in the latter denominator. These
single-trajectory readouts do not receive the final procedure's additional proposals or selection.
The recovered raw endpoint gives 323/1,408 exact first-draw outputs (22.9\%) and 1,658/7,040
across all five draws (23.6\%). All 880 conditioned replay batches reproduce every saved prediction
tensor exactly. The independently trained shared-null comparison remains unmeasured.

The selected-product inverse search returns a candidate for 1,392/1,408 requests (98.9\%). Of
1,456 executor-specific candidate checks, 1,387 pass the admitted replay/source checks (95.3\%).
Sixteen requests return no candidate and five return candidates without an exact accepted replay.
Sixty-four of the 69 rejected tuples are alternative executor bindings in the thiol-yne family,
where all 64 selected products still have an accepted exact decomposition.
All inverse searches finish without recorded search-bound flags; no request has multiple distinct
accepted constitutional role/component tuples among the returned candidates. This is bounded
computational acceptance, not experimentally adjudicated chemical precision or exhaustive route
uniqueness. The matched three-family comparison remains
\PendingValue{common-program comparison}.

\begin{table}[H]
\caption{\textbf{Exact assembly across 22 ionizable-lipid reaction families.}
Exact-L1 percentages for one checkpoint and 64 TRAIN-derived development requests per family.
Raw, Terminal and Cyclic raw use the first draw. Raw retains the requested core and layout;
Cyclic changes only the program ID. Terminal is the final argmax. Final additionally uses five draws,
checked assembly, proposals and selection, with a larger candidate budget than the paired readouts.
Decomposition coverage and candidate acceptance refer to selected Final products; their denominators
are requests and executor-specific candidate checks, respectively. Dashes denote unavailable
measurements, not zero. Independent-fit uncertainty is not estimated.}
\label{tab:22-l1}
\centering\footnotesize\setlength{\tabcolsep}{3pt}\renewcommand{\arraystretch}{1.12}
\begin{tabular}{@{}>{\raggedright\arraybackslash}p{.28\textwidth}cccccc@{}}
\toprule
Family & \shortstack{Raw\\L1 (\%)} & \shortstack{Terminal\\L1 (\%)} & \shortstack{Final\\L1 (\%)} & \shortstack{Shared\\null (\%)} & \shortstack{Cyclic raw\\L1 (\%)} & \shortstack{Decomp. cov. /\\cand. accept. (\%)} \\
\midrule
A3 amine--aldehyde--alkyne & 20.3 & 21.9 & 95.3 & \PendingCell & 29.7 & 96.9/98.4 \\
Acid--epoxide diester & 25.0 & 45.3 & 98.4 & \PendingCell & 25.0 & 98.4/100.0 \\
AEMA & 21.9 & 26.6 & 98.4 & \PendingCell & 21.9 & 98.4/100.0 \\
Aldehyde Ugi 3-CR & 65.6 & 71.9 & 98.4 & \PendingCell & 70.3 & 100.0/98.4 \\
Aldehyde Ugi 4-CR & 48.4 & 54.7 & 100.0 & \PendingCell & 48.4 & 100.0/100.0 \\
$\alpha$-Isocyanoester dihydroimidazole & 26.6 & 29.7 & 100.0 & \PendingCell & 10.9 & 100.0/100.0 \\
Amine alkylation & 0.0 & 0.0 & 96.9 & \PendingCell & 0.0 & 96.9/100.0 \\
Amine--epoxide & 0.0 & 1.6 & 100.0 & \PendingCell & 1.6 & 100.0/100.0 \\
Aryl reductive amination & 1.6 & 1.6 & 90.6 & \PendingCell & 0.0 & 95.3/95.1 \\
Aza-Michael acrylamide & 7.8 & 9.4 & 100.0 & \PendingCell & 9.4 & 100.0/100.0 \\
Aza-Michael acrylate & 1.6 & 3.1 & 100.0 & \PendingCell & 6.2 & 100.0/100.0 \\
Disulfide Michael & 0.0 & 0.0 & 98.4 & \PendingCell & 0.0 & 98.4/100.0 \\
Epoxide O-acylation & 4.7 & 9.4 & 98.4 & \PendingCell & 3.1 & 98.4/100.0 \\
iPhos & 57.8 & 60.9 & 100.0 & \PendingCell & 64.1 & 100.0/100.0 \\
Ketone--isocyanide amide & 51.6 & 57.8 & 100.0 & \PendingCell & 37.5 & 100.0/100.0 \\
Ketone Ugi 4-CR & 4.7 & 3.1 & 96.9 & \PendingCell & 1.6 & 96.9/100.0 \\
Maleate & 0.0 & 0.0 & 100.0 & \PendingCell & 1.6 & 100.0/100.0 \\
O-esterification & 39.1 & 40.6 & 100.0 & \PendingCell & 35.9 & 100.0/100.0 \\
Passerini & 68.8 & 78.1 & 100.0 & \PendingCell & 54.7 & 100.0/100.0 \\
Preassembled thiol-yne & 0.0 & 3.1 & 100.0 & \PendingCell & 3.1 & 100.0/50.0 \\
Reductive amination & 0.0 & 1.6 & 95.3 & \PendingCell & 0.0 & 95.3/100.0 \\
Thiolactone & 59.4 & 71.9 & 100.0 & \PendingCell & 64.1 & 100.0/100.0 \\
\midrule
Equal-family macro & 22.9 & 26.9 & 98.5 & \PendingCell & 22.2 & 98.9/97.4 \\
Pooled all requests & 22.9 & 26.9 & 98.5 & \PendingCell & 22.2 & 98.9/95.3 \\
Worst-family result & 0.0 & 0.0 & 90.6 & \PendingCell & 0.0 & 95.3/50.0 \\
\bottomrule
\end{tabular}
\par\smallskip\raggedright\scriptsize
Worst-family entries are separate column-wise minima and need not describe the same family.
Candidate checks include rejected returned tuples; earlier inverse-template pruning is not fully
logged. Acceptance is not precision over all raw inverse outcomes or an independent chemistry audit.
\end{table}

\subsection{Program conditioning and training diagnostics}
On the same 1,408-request first draw, cycling the program ID reduces raw exact L1 from 323 to
313 (88 gains and 98 losses), while valid connected outputs increase from 866 to 927. Terminal
argmax gives 379 versus 368 exact outputs and 996 versus 1,067 valid connected outputs. Across
all five draws (7,040 trajectories per arm), raw exact counts are 1,658 versus 1,521 and terminal
exact counts are 1,910 versus 1,794. These are repeated trajectories from one fit and the same
requests, not independent training seeds. The intervention retains the true core, roles, layout
and depth; the small first-draw exact difference and opposing validity change limit its interpretation.

On the 352-request paired population, changing only the program identity at inference reduces D1
exact L1 from 140/352 to 105/352, a 9.94 percentage-point decrease. Raw endpoint exact counts are
80 versus 71, and terminal-argmax counts are 91 versus 85. The cyclic intervention retains the true
reaction core, roles, layout and depth. It therefore measures sensitivity to the program identity
within otherwise correct structural context, not the full effect of synthesis information.
The independently trained shared-null contrast remains \PendingValue{trained null comparison}.

Additional training did not improve D1 exact yield. With another 31,680 presentations per family,
220 large-batch updates and 2,640 small-batch updates produced 135 and 122 exact outputs,
respectively, versus 140 for their common initializer (all /352). Hardware, computation and
PCGrad grouping differed, so matched exposure does not imply matched cost. A second experiment
compared two 66-update continuations, each adding 9,504 presentations per family. Ordinary
categorical and coarse parent-group training gave 136 and 135 exact outputs, respectively.
Although coarse training reduced nonconsecutive-parent CAL errors from 1,164 to 1,026 of 7,936,
it did not preserve the original generation floors. Neither continuation replaced the original
checkpoint (Figure~\ref{fig:22-training}; Table~\ref{tab:22-architecture}).

\subsection{Decoder contributions and limitations}
On the 1,408-request cohort's first draw, the saved D0 constrained readout and D1 completion
give 554 and 576 exact outputs, respectively (39.3\% and 40.9\%), compared with 379 (26.9\%)
for terminal argmax. D0 includes core reservation and is not a raw flow endpoint. The selected
1,387/1,408 result additionally uses five draws, construction proposals and selection, so this
larger gain cannot be attributed to a single decoder rule or increased model learning.

For the original checkpoint on the 352-request diagnostic, raw endpoint, terminal-argmax and D1
readouts yield 80, 91 and 140 exact products, respectively. A fixed fallback among these already
computed candidates increases D1 exact yield to 145/352 and limited-design yield from 111 to
116/352, without new model or chemistry calls. This is a selection gain. In a separate paired
192-request repair diagnostic covering A3 and the two Ugi 4-CR families, exact L1 remains 187/192
while limited-design passes increase from 140 to 146. Additional attachment proposals change three
structures but add no design passes. These repairs were not substituted into the 1,408-request
routing population. Table~\ref{tab:22-decoder} retains the separate denominators and adverse results.
Matched whole-graph versus factorized capacity, cross-attention and gradient-control contrasts
remain \PendingValue{architecture comparisons}.

\subsection{Precursor novelty and diversity}
The selected population contains 1,407 distinct products among 1,408 requests. Product novelty
relative to TRAIN occurs in 1,216 observations (1,215 distinct identities). Exact products with at
least one globally TRAIN-novel recovered component occur in 1,050/1,408 requests (74.6\%),
representing 1,049 distinct products. This gives 745.0 distinct component-novel products per 1,000
requests. Role-relative novelty occurs in 1,081 observations (1,080 distinct products).
Component counts remain uneven: the pooled isocyanide role contains 29 identities among 190
observations, with Shannon and inverse-Simpson effective counts of 10.21 and 6.09
(Table~\ref{tab:22-diversity}). High product uniqueness therefore does not ensure broad component use.

TRAIN-relative novelty does not establish catalogue escape, recovery of deliberately withheld
components or source-disjoint generalization. Only 15/2,059 guarded CAL molecules are wholly
component-disjoint. The newer 1,859-molecule exact-source and representation-qualified CAL union
contains 12 such cases, all in two aza-Michael families; it differs from the older 1,821-reference
distribution cohort below. Final held-component and source-disjoint generation results remain
\PendingValue{generalization evaluation} (Table~\ref{tab:22-generalization}).

\subsection{Structural quality and agreement with lipid references}
All 1,408 selected structures are valid and connected, but only 1,324/1,408 (94.0\%) satisfy exact
L1 together with all applicable limited design checks (Table~\ref{tab:22-quality-metrics}).
There are 69 failures of the combined ring predicate. The scoped retained-head predicate fails in
5/64 aldehyde Ugi 4-CR requests; the other 1,344 structures are outside its qualified scope and
are not counted as passes. Chemistry-context review is flagged for 41/1,408 structures (2.9\%).
The overlapping flags include hydrated or hemiacetal carbons (24), peroxide groups (8), carbons
bearing at least three oxygens (6), other allenes (5), a small unsaturated ring (1) and a heteroatom
cumulene (1). Limited design without context flags holds for 1,292/1,408 (91.8\%).

\begin{table}[H]
\caption{\textbf{Structural quality of the selected development population.}
Counts retain all 1,408 requests except explicitly scoped head and ring predicates. The limited-design
conjunction includes exact L1 and the qualified chemical, head and ring checks; it is not a broad
empirical lipid-realism assessment. TRAIN support measures observed structural features, not external validation.}
\label{tab:22-quality-metrics}\centering\footnotesize
\setlength{\tabcolsep}{3pt}\renewcommand{\arraystretch}{1.12}
\begin{tabular}{@{}>{\raggedright\arraybackslash}p{.61\textwidth}rc@{}}\toprule
Metric & Selected count & (\%)\\\midrule
Valid molecules / requests & 1,408/1,408 & 100.0 \\
Connected molecules / requests & 1,408/1,408 & 100.0 \\
Exact L1 + qualified limited design / requests & 1,324/1,408 & 94.0 \\
Limited design without context flags / requests & 1,292/1,408 & 91.8 \\
Justified narrow chemical alerts / requests & 0/1,408 & 0.0 \\
Chemistry-context flags / requests & 41/1,408 & 2.9 \\
Required head retained / scoped requests & 59/64 & 92.2 \\
Combined ring predicate / requests & 1,339/1,408 & 95.1 \\
Role-local cycle allocation / requests & 1,404/1,408 & 99.7 \\
Requested ring-size agreement / applicable requests & 720/789 & 91.3 \\
Mean per-product TRAIN atom-environment coverage & 1,408 products & 98.5 \\
Ring-system TRAIN support / ring-containing products & 589/810 & 72.7 \\
Exact L1 + product/component features observed in TRAIN & 974/1,408 & 69.2 \\
\bottomrule\end{tabular}
\par\smallskip\raggedright\scriptsize
Ring-size checks are inapplicable to 619 requests. For ring-system TRAIN support, 221 ring-containing
products have unknown support and 598 are ring-free; neither group is silently counted as a pass.
The atom-weighted coverage is 77,273/78,532 (98.4\%), distinct from the per-product mean.
Source-positive controls pass the narrow chemical-alert screen in 29/29 cases, but cover only two families
and two studies, with zero individually characterized products in that control set. The other 20
families lack these controls. Absence of an alert does not establish stability, pKa or delivery activity.
Definitions and family-level counts: Tables~\ref{tab:22-structure}--\ref{tab:22-metric-definitions}.
\end{table}

The distribution comparison uses 1,821 exact-source-qualified \emph{synthetic CAL structures};
these are not an independent empirical lipid reference. Equal-family fingerprint precision is
56.9\%, coverage is 49.1\%, and mean within-family ECFP4 diversity is 0.552
(Table~\ref{tab:22-realism-metrics}). Precision varies from 1/64 for ketone Ugi 4-CR to 64/64 for
aza-Michael acrylate. Aryl reductive amination has 2/63 supported distinct outputs. These failures
remain visible despite high exact L1. Compared with the raw first draw, the full selected procedure
increases distinct fingerprint-supported output yield from 399/1,408 to 801/1,408, while mean
within-family diversity decreases from 0.585 to 0.552. This comparison combines decoding, additional
proposals and selection; it does not isolate a learned-model improvement. A3 has 36/64 supported outputs; its normalized Wasserstein
distances for cLogP and rotatable bonds are 1.133 and 1.658. Ketone Ugi 4-CR has an sp$^3$-carbon
fraction distance of 1.082. All 24 descriptor panels are reported by family
(Table~\ref{tab:22-descriptors}). Independent empirical-reference comparisons, grouped two-sample
classification and blinded structural review remain unmeasured.

\begin{table}[H]
\caption{\textbf{Distributional agreement with synthetic CAL references.}
Each arm has 1,408 requests, 64 per family. Raw and Terminal use the first conditioned draw;
Selected additionally uses five draws, checked construction and selection. These stages have
unequal candidate budgets. Values are equal-family means of 22 within-family statistics;
precision and diversity are conditional on connected molecules (fingerprints also deduplicate).
All-request support yields retain invalid outputs. TRAIN controls are a reference diagnostic,
not independent quality evidence.}
\label{tab:22-realism-metrics}\centering\footnotesize
\setlength{\tabcolsep}{3pt}\renewcommand{\arraystretch}{1.12}
\begin{tabular}{@{}>{\raggedright\arraybackslash}p{.49\textwidth}cccc@{}}\toprule
Metric & Raw & Terminal & Selected & TRAIN\\\midrule
Connected / all requests (\%) & 61.5 & 70.7 & 100.0 & 100.0 \\
Fingerprint precision (\%) & 41.6 & 47.6 & 56.9 & 75.1 \\
Fingerprint coverage (\%) & 34.1 & 36.4 & 49.1 & 59.0 \\
Fingerprint recall (\%) & 81.2 & 80.0 & 82.1 & 88.4 \\
Distinct fingerprint-supported products / requests (\%) & 28.3 & 34.7 & 56.9 & 75.1 \\
Nearest-reference Tanimoto similarity & 0.579 & 0.604 & 0.638 & 0.693 \\
Within-family ECFP4 diversity & 0.585 & 0.559 & 0.552 & 0.529 \\
Synthetic-CAL descriptor precision (\%) & 57.4 & 58.6 & 59.4 & 69.7 \\
Synthetic-CAL descriptor coverage (\%) & 32.0 & 33.1 & 39.7 & 44.0 \\
Descriptor-supported observations / requests (\%) & 36.6 & 42.2 & 59.4 & 69.7 \\
Descriptor distances and dispersion (24 features, per family) & \multicolumn{4}{c}{Table~\ref{tab:22-descriptors}} \\
Independent descriptor precision / coverage & \PendingCell & \PendingCell & \PendingCell & \PendingCell \\
Independent source-grouped classifier AUC & \PendingCell & \PendingCell & \PendingCell & \PendingCell \\
\bottomrule\end{tabular}
\par\smallskip\raggedright\scriptsize
Connected counts are 866/1,408 (Raw), 996/1,408 (Terminal) and 1,408/1,408 (Selected).
ECFP4 uses Morgan radius 2 and 2,048 bits; neighborhoods use up to five neighbors.
Descriptor neighborhoods use 24 TRAIN-scaled descriptors and the same synthetic reference.
Coverage is the equal-family mean of covered reference neighborhoods, not the pooled ratio.
Similarity to this reference is not a probability of chemical realism. The predecessor selection
and matched cyclic readouts remain in Table~\ref{tab:22-realism}.
\end{table}

\subsection{Recursive computational makeability and strict dossiers}
The primary route criterion requires every assembly component to be directly vendor-listed or to
have a complete computational path to vendor-listed terminal materials. Planner-generated paths
and forward-applied published transforms may contribute without exact-substrate experimental
precedent. On the unchanged 1,408-product population, 616/1,408 (43.75\%) meet this criterion,
including 606 satisfying limited design and 604 also free of context flags
(Table~\ref{tab:22-routes}). L2 readiness is 1,189/1,408 (84.4\%); it does not require terminal
listing closure. Direct-only L3 covers 119/1,408 (8.5\%). Strict exact-source/current-item dossiers
remain a separate secondary outcome, closing 27/1,408 requests (1.9\%).

Makeability ranges from 0/64 for reductive amination and 2/64 for ketone Ugi 4-CR to 58/64 for
ketone--isocyanide amide. The complete denominator includes 3,723 required component branches.
The latest listing evidence is dated September 25, 2026 (local date), with a 30-day policy window.
A vendor-directory listing does not establish stock, price or an experimentally executed route.
The primary count rose from 581 to 616 through additional evidence, without changing molecules,
L1 checks or structural predicates. Missing listings and bounded route searches remain unresolved,
not evidence of unsynthesizability. Neither outcome establishes experimental synthesis or delivery.
